\section{代数叠的局部结构}\label{src:7.7}
\begin{epigraph}
因此, 人类的一切知识始于直观, 由此进到概念, 最后止于理念.
\epigraphauthor{Immanuel Kant}
\end{epigraph}
我们证明在线性约化稳定子点附近的代数叠局部结构定理~\bookref{7.7.1}. 从思想上说, 形如 $[\Spec A/G]$ 且 $G$ 线性约化的商叠, 是代数叠在线性约化稳定子点附近的基本构件, 正如仿射概形是概形和代数空间的基本构件. 这使许多代数叠的性质可归约到已充分理解的 $[\Spec A/G]$ 情形. 此定理推广了 Deligne--Mumford 叠局部结构定理~\bookref{5.3.1}, 并有类似应用. 例如, 我们将用它构造良模空间 (定理~\bookref{7.10.1}), 类似于在 Keel--Mori 定理~\bookref{5.4.6} 的证明中用 Deligne--Mumford 情形构造粗模空间.
\begin{theorem}[代数叠的局部结构定理]\label{src:7.7.1}
设 $\mathcal X$ 是代数闭域 $\Bbbk$ 上有限型且具有仿射对角态射的代数叠. 对每个具有线性约化稳定子 $G_x$ 的点 $x\in\mathcal X(\Bbbk)$, 存在仿射平展态射
\begin{align*}
f:([\Spec A/G_x],w)\longrightarrow(\mathcal X,x),
\end{align*}
它在 $w$ 处诱导稳定子群的同构.
\end{theorem}
\begin{remark}\label{src:7.7.2}
若 $x\in|\mathcal X|$ 是光滑点, 还可得到平展态射 $([\Spec A/G_x],w)\to([T_{\mathcal X,x}/G_x],0)$, 其中 $T_{\mathcal X,x}$ 是带有 $G_x$-表示结构的 Zariski 切空间. 这一补充既由证明得到, 也可通过对 $[\Spec A/G_x]$ 应用 Luna 平展切片定理~\bookref{7.7.5} 得到. 因而 $\mathcal X$ 的平展局部性质可归约为 $T_{\mathcal X,x}$ 的 $G_x$-等变性质; 对模问题, 这就转化为研究一阶变形空间作为自同构群表示的性质.
\end{remark}
结合此定理与 Luna 基本引理~\bookref{7.5.30}, 得到:
\begin{corollary}[良模空间的局部结构]\label{src:7.7.3}
设 $\mathcal X$ 是代数闭域 $\Bbbk$ 上有限型且具有仿射对角态射的代数叠, 并存在良模空间 $\pi:\mathcal X\to X$. 对每个闭点 $x\in|\mathcal X|$, 存在 $\pi(x)$ 的平展邻域 $W\to X$ 与笛卡尔图
\begin{align*}
\begin{tikzcd}[ampersand replacement=\&]
{[\Spec A/G_x]}\arrow[r]\arrow[d]\&\mathcal X\arrow[d,"{\pi}"]\\
W=\Spec A^{G_x}\arrow[r]\&X.
\end{tikzcd}
\end{align*}
\end{corollary}
一个好处是可给出良模空间的等价而更直观的定义.
\begin{corollary}[良模空间的刻画]\label{src:7.7.4}
设 $\mathcal X$ 是代数闭域 $\Bbbk$ 上有限型且具有仿射对角态射的代数叠. 到代数空间的态射 $\pi:\mathcal X\to X$ 是良模空间, 当且仅当每个 $x\in X$ 有平展邻域 $\Spec B\to X$, 使得 $\mathcal X\times_X\Spec B\cong[\Spec A/G]$, 其中 $G$ 线性约化且 $B=A^G$.
\end{corollary}
\emph{本节安排.} 先讨论 Luna 平展切片定理~\bookref{7.7.5}, 它在已知 $\mathcal X=[\Spec B/G]$ 且 $G$ 线性约化时, 用一个漂亮论证显式构造平展邻域. 局部结构定理~\bookref{7.7.1} 的证明远不如此显式, 需要 (1) 变形理论, (2) 凝聚完备性, (3) 凝聚 Tannaka 对偶, (4) Artin 逼近或等变 Artin 代数化.

令 $\mathcal T=[T_{\mathcal X,x}/G_x]$, 变形理论产生 $x$ 与 $0$ 的第 $n$ 阶幂零加厚的嵌入 $\mathcal X_n\to\mathcal T_n$. 关键是证明闭态射系统 $(\mathcal X_n\to\mathcal T_n)$ 可代数化. 第一步是有效化: 令 $\pi:\mathcal T\to T:=T_{\mathcal X,x}//G_x$, 则 $\widehat{\mathcal T}:=\mathcal T\times_T\Spec\widehat{\mathcal O}_{T,\pi(0)}$ 凝聚完备 (定理~\bookref{7.6.13}). 因而可构造延拓 $\mathcal X_n\hookrightarrow\mathcal T$ 的闭子叠 $\widehat{\mathcal X}\hookrightarrow\widehat{\mathcal T}$, 再用凝聚 Tannaka 对偶~\bookref{7.6.9} 构造延拓 $\mathcal X_n\to\mathcal X$ 的态射 $\widehat{\mathcal X}\to\mathcal X$. 若 $x$ 光滑, 在 GIT 商 $T_{\mathcal X,x}//G_x$ 上的 Artin 逼近产生平展邻域 $U\to T_{\mathcal X,x}//G_x$, 使 $\pi^{-1}(U)\to\mathcal X$ 代数化 $\widehat{\mathcal T}\to\mathcal X$. 一般情形更复杂, 需要等变 Artin 代数化 (定理~\bookref{7.7.18}).

\subsection{Luna 平展切片定理}\label{src:7.7.1-subsection}
局部结构定理~\bookref{7.7.1} 来自等变几何中的 Luna 平展切片定理.
\begin{theorem}[Luna 平展切片定理]\label{src:7.7.5}
设 $G$ 是代数闭域 $\Bbbk$ 上的线性约化群, $X$ 是带 $G$-作用的有限型仿射 $\Bbbk$-概形. 若 $x\in X(\Bbbk)$ 的稳定子线性约化, 则存在 $G_x$-不变, 局部闭且仿射的子概形 $W\subseteq X$, 使诱导映射
\begin{align}\label{src:7.7.6}
[W/G_x]\longrightarrow[X/G]
\end{align}
仿射平展. 若轨道 $Gx\subseteq X$ 还闭, 则有笛卡尔图
\begin{align*}
\begin{tikzcd}[ampersand replacement=\&]
{[W/G_x]}\arrow[r]\arrow[d]\&{[X/G]}\arrow[d]\\
W//G_x\arrow[r]\&X//G,
\end{tikzcd}
\end{align*}
其中 $W//G_x\to X//G$ 也平展. 此外, 若 $x$ 是光滑点, 记轨道的法空间 $N_x=T_{X,x}/T_{Gx,x}$, 可安排有笛卡尔图
\begin{align*}
\begin{tikzcd}[ampersand replacement=\&]
{[N_x/G_x]}\arrow[d]\&{[W/G_x]}\arrow[l]\arrow[r]\arrow[d]\&{[X/G]}\arrow[d]\\
N_x//G_x\&W//G_x\arrow[l]\arrow[r]\&X//G,
\end{tikzcd}
\end{align*}
所有横箭头平展, 且 $[W/G_x]\to[N_x/G_x]$ 来自 $G_x$-不变平展态射 $W\to N_x$.
\end{theorem}
\begin{remark}\label{src:7.7.7}
若轨道 $G_x$ 闭, 则 Matsushima 定理~\bookref{7.5.23} 说明稳定子 $G_x$ 线性约化.
\end{remark}
\begin{remark}\label{src:7.7.8}
也可作 $G$-等变表述: $G$ 自然作用于商 $G\times^{G_x}W:=(G\times W)/G_x$, 且有等同 $[W/G_x]\cong[(G\times^{G_x}W)/G]$, 同样有 $W//G_x\cong(G\times^{G_x}W)/G$ (见习题~\bookref{4.4.19}). 态射 \bookeqref{7.7.6} 对应于 $G$-等变平展态射 $G\times^{G_x}W\to X$.
\end{remark}
证明依赖于一个 $G_x$-不变态射 $X\to T_{X,x}$ 的存在性, 称为 Luna 映射. 这里略滥用记号 $T_{X,x}$: 它是 $\Bbbk$-向量空间, 通过 $\Spec(\operatorname{Sym}^*T_{X,x}^\vee)$ 将其视为概形.
\begin{lemma}[Luna 映射]\label{src:7.7.9}
设 $G$ 是代数闭域 $\Bbbk$ 上的线性约化群, $X$ 是带 $G$-作用的有限型仿射 $\Bbbk$-概形. 若 $x\in X(\Bbbk)$ 的稳定子线性约化, 则存在将 $x$ 送到原点的 $G_x$-等变态射
\begin{align}\label{src:7.7.10}
f:X\longrightarrow T_{X,x}.
\end{align}
若 $X$ 在 $x$ 处光滑, 则 $f$ 在 $x$ 处平展.
\end{lemma}
\begin{proof}
写 $X=\Spec A$, 令 $\mathfrak m\subseteq A$ 为 $x$ 的极大理想. 则 $\mathfrak m$ 与 $\mathfrak m/\mathfrak m^2$ 是 $G_x$-表示, 所以 $G_x$ 自然作用于切空间 $T_{X,x}:=\Spec(\operatorname{Sym}^*(\mathfrak m/\mathfrak m^2))$. 因 $G_x$ 线性约化, $G_x$-表示的满射 $\mathfrak m\to\mathfrak m/\mathfrak m^2$ 有截面 $\mathfrak m/\mathfrak m^2\hookrightarrow\mathfrak m$. 它诱导等变环映射 $\operatorname{Sym}^*(\mathfrak m/\mathfrak m^2)\to A$, 从而诱导将 $x$ 送到原点的 $G_x$-等变态射 $f:\Spec A\to T_{X,x}$. 若 $x$ 光滑, 因 $f$ 在 $x$ 处诱导切空间同构, 定理~\bookref{A.3.2} 说明它在该处平展.
\end{proof}
\begin{proof}[定理~\bookref{7.7.5} 的证明]
因 $X$ 仿射有限型, 可取有限维 $G$-表示 $V$ 与 $G$-等变闭浸入 $X\hookrightarrow\mathbb A(V)$ (命题~\bookref{B.1.21}). 若 $W\subseteq\mathbb A(V)$ 是仿射 $G_x$-不变局部闭子概形, 且 $[W/G_x]\to[\mathbb A(V)/G]$ 平展, 则 $W':=W\cap X$ 及 $[W'/G_x]\to[X/G]$ 也如此. 因此立即归约到 $x$ 光滑的情形. 此时 Luna 映射 $f:X\to T_{X,x}$ 是 $G_x$-不变的, 在 $x$ 处平展且 $f(x)=0$ (引理~\bookref{7.7.9}). 子空间 $T_{Gx,x}\subseteq T_{X,x}$ 为 $G_x$-不变; 再用 $G_x$ 线性约化, 满射 $T_{X,x}\to N_x=T_{X,x}/T_{Gx,x}$ 有截面 $N_x\hookrightarrow T_{X,x}$. 定义 $W$ 为 $N_x$ 在 $f$ 下的原像:
\begin{align*}
\begin{tikzcd}[ampersand replacement=\&]
W\arrow[r]\arrow[d]\&N_x\arrow[d]\\
X\arrow[r,"{f}"']\&T_{X,x}.
\end{tikzcd}
\end{align*}
映射 $f:[W/G_x]\to[X/G]$ 与 $g:[W/G_x]\to[N_x/G_x]$ 在 $w$ 处诱导切空间与稳定子群的同构, 所以它们都在 $x\in W$ 处平展; 等价地, $G$-等变映射 $G\times^{G_x}W\to X$ 与 $G\times^{G_x}W\to G\times^{G_x}N_x$ 在 $(\operatorname{id},x)$ 处平展. 有交换图
\begin{align*}
\begin{tikzcd}[ampersand replacement=\&]
{[N_x/G_x]}\arrow[d]\&{[W/G_x]}\arrow[l,"{g}"']\arrow[r,"{f}"]\arrow[d]\&{[X/G]}\arrow[d]\\
N_x//G_x\&W//G_x\arrow[l]\arrow[r]\&X//G_x.
\end{tikzcd}
\end{align*}
其中 $f,g$ 在 $x$ 处平展, 在该处保持稳定子群, 并将 $x$ 映到闭点. 应用 Luna 基本引理~\bookref{7.5.30}, 可将 $W$ 换为 $x$ 的 $G_x$-等变仿射开邻域, 使各方块笛卡尔. 另见 \cite[第 97 页]{Lun73} 或 \cite[定理 5.3]{Dré04}.
\end{proof}
若已知 $\mathcal X$ 是正规拟射影概形的商叠, 局部结构定理可由直接论证得到. 此情形足以处理许多模问题, 如特征 $0$ 下的 $\mathcal B\!un_{r,d}^{\mathrm{ss}}(C)$.
\begin{exercise}\label{src:7.7.11}
若代数闭域 $\Bbbk$ 上的连通仿射代数群 $G$ 作用于正规有限型 $\Bbbk$-概形 $X$, 且 $x\in X(\Bbbk)$ 的稳定子线性约化, 证明存在 $G_x$-不变, 局部闭且仿射的子概形 $W\hookrightarrow X$, 使 $[W/G_x]\to[X/G]$ 平展.

提示: 用 Sumihiro 线性化定理~\bookref{B.1.33} 归约到 $X=\mathbb P(V)$. 取在 $x$ 处非零的齐次多项式 $f$, 使 $\mathbb P(V)_f$ 为 $G_x$-不变, 再考虑 $G_x$-等变平展映射 $\mathbb P(V)_f\to T_x\mathbb P(V)$, 如 Luna 平展切片定理的证明.
\end{exercise}

\subsection{变形理论}\label{src:7.7.2-subsection}
局部结构定理~\bookref{7.7.1} 的证明需要一些代数叠的变形理论.
\begin{proposition}\label{src:7.7.12}
考虑具有仿射对角态射的 Noether 代数叠的交换图
\begin{align*}
\begin{tikzcd}[ampersand replacement=\&]
\mathcal W\arrow[r,"{f}"]\arrow[d,hook]\&\mathcal X\arrow[d]\\
\mathcal W'\arrow[r]\arrow[ur,dashed]\&\mathcal Y,
\end{tikzcd}
\end{align*}
其中 $\mathcal X\to\mathcal Y$ 光滑仿射, $\mathcal W\hookrightarrow\mathcal W'$ 是由平方零理想层定义的闭浸入. 若 $\mathcal W$ 上同调仿射, 则该图有提升.
\end{proposition}
\begin{proof}
$\mathcal W$ 仿射时结论来自无穷小提升判据~\bookref{A.3.1}. 为归约到这一情形, 取仿射概形的光滑展示 $U'\to\mathcal W'$, 定义 $U:=U'\times_{\mathcal W'}\mathcal W$. 因 $\mathcal W$ 的对角态射仿射, 每个 $n$ 重纤维积 $(U/\mathcal W)^n:=U\times_{\mathcal W}\cdots\times_{\mathcal W}U$ 仿射. 有交换图
\begin{align*}
\begin{tikzcd}[ampersand replacement=\&,column sep=small]
(U/\mathcal W)^2\arrow[r,shift left]\arrow[r,shift right]\arrow[d,hook]\&U\arrow[r,"{q_1}"]\arrow[d,hook]\&\mathcal W\arrow[r,"{f}"]\arrow[d,hook]\&\mathcal X\arrow[d]\\
(U'/\mathcal W')^2\arrow[r,shift left,"{p_1}"]\arrow[r,shift right,"{p_2}"']\&U'\arrow[r]\arrow[urr,"{f'_U}",near end]\&\mathcal W'\arrow[r]\arrow[ur,dashed,"{f'}"]\&\mathcal Y,
\end{tikzcd}
\end{align*}
其中已取提升 $f'_U:U'\to\mathcal X$. 在 $\mathcal W$ 上定义凝聚层 $F=f^*(\Omega_{\mathcal X/\mathcal Y}^\vee)\otimes J$, $J$ 为定义闭浸入的理想层. 习题~\bookref{7.1.16} 说明提升 $U'\to\mathcal X$ 的集合是 $\Gamma(U,q_1^*F)$ 下的挠子, 即任一提升与 $f'_U$ 的差为其中一个元素. 因 $\mathcal X\to\mathcal Y$ 可表, 为使 $f'_U$ 下降为 $f':\mathcal W'\to\mathcal X$, 需安排 $f'_U\circ p_1=f'_U\circ p_2$. 令 $q_n:(U/\mathcal W)^n\to\mathcal W$. 差 $f'_U\circ p_1-f'_U\circ p_2$ 可视为 $\Gamma((U/\mathcal W)^2,q_2^*F)$ 中的元素. 因 $q_1$ 满射, 光滑且仿射, 有拟凝聚层正合列
\begin{align*}
0\to F\to q_{1,*}q_1^*F\to q_{2,*}q_2^*F\to q_{3,*}q_3^*F\to\cdots;
\end{align*}
见习题~\bookref{3.1.3}. $\mathcal W$ 上同调仿射, 故取整体截面得到正合列
\begin{align*}
\Gamma(U,q_1^*F)\xrightarrow{d_0}\Gamma((U/\mathcal W)^2,q_2^*F)\xrightarrow{d_1}\Gamma((U/\mathcal W)^3,q_3^*F),\\
s\longmapsto p_1^*s-p_2^*s,\qquad t\longmapsto p_{12}^*t-p_{13}^*t+p_{23}^*t.
\end{align*}
可检验 $d_1(f'_U\circ p_1-f'_U\circ p_2)=0$, 故有 $s\in\Gamma(U,q_1^*F)$ 满足 $d_0(s)=f'_U\circ p_1-f'_U\circ p_2$. 用 $s$ 修改提升后, 差变为 $0$, 所以它下降为 $f'$.
\end{proof}
\begin{remark}\label{src:7.7.13}
也可证明此变形问题的障碍属于 $\operatorname{Ext}^1_{\mathcal O_{\mathcal W}}(f^*\Omega_{\mathcal X/\mathcal Y},J)=H^1(\mathcal W,f^*(\Omega_{\mathcal X/\mathcal Y}^\vee)\otimes J)$, 因 $\mathcal W$ 上同调仿射而消失. 更一般的结果见 \cite[定理 1.5]{Ols06}.
\end{remark}
\begin{exercise}\label{src:7.7.14}
设 $\mathcal W\hookrightarrow\mathcal W'$ 是有限型 $\Bbbk$-代数叠之间的闭浸入, 两者有仿射对角态射, 该浸入由平方零理想层定义. 设 $G$ 是仿射 $\Bbbk$-代数群. 若 $\mathcal W$ 上同调仿射, 证明每个主 $G$-丛 $\mathcal P\to\mathcal W$ 延拓为主 $G$-丛 $\mathcal P'\to\mathcal W'$.

提示: 用光滑下降以及概形上主 $G$-丛的变形理论 (习题~\bookref{C.2.6}).
\end{exercise}
\begin{remark}\label{src:7.7.15}
习题~\bookref{7.7.14} 等价于将 $f:\mathcal W\to BG$ 变形为 $\mathcal W'\to BG$, 类似命题~\bookref{7.7.12}, 但 $\mathcal X=BG\to\mathcal Y=\Spec\Bbbk$ 非仿射. 到 $BG$ 的映射对应于主 $G$-丛, 故问题等价于将 $P\to\mathcal W$ 变形为 $P'\to\mathcal W'$. 与 $\mathcal W$ 为概形的情形 (习题~\bookref{C.2.6}) 一样, 障碍属于 $H^2(\mathcal W,\mathfrak g\otimes J)$, $J$ 为定义 $\mathcal W\hookleftarrow\mathcal W'$ 的理想层. 余切复形的一般理论 (见注记~\bookref{C.3.7}, \bookref{C.7.6}) 说明障碍属于 $\operatorname{Ext}^1_{\mathcal O_{\mathcal W}}(Lf^*L_{BG/\Bbbk},J)$. 对复合 $\Spec\Bbbk\xrightarrow pBG\to\Spec\Bbbk$, 有正合三角 $p^*L_{BG/\Bbbk}\to L_{\Bbbk/\Bbbk}\to L_{\Bbbk/BG}$. 因 $L_{\Bbbk/\Bbbk}=0$, 得到 $p^*L_{BG/\Bbbk}=L_{\Bbbk/BG}[-1]\cong\mathfrak g^\vee[-1]$, 以及 $L_{BG/\Bbbk}\cong\mathfrak g^\vee[-1]$, 其中 Lie 代数 $\mathfrak g$ 配备伴随表示. 因此
\begin{align*}
\operatorname{Ext}^1_{\mathcal O_{\mathcal W}}(Lf^*L_{BG/\Bbbk},J)=H^1(\mathcal W,f^*\mathfrak g[1]\otimes J)=H^2(\mathcal W,\mathfrak g\otimes J).
\end{align*}
因 $\mathcal W$ 上同调仿射且对角态射仿射, 此上同调群为 $0$, 障碍消失.
\end{remark}
还需要以下闭浸入与同构的判据.
\begin{lemma}\label{src:7.7.16}
设 $f:\mathcal X\to\mathcal Y$ 是代数闭域 $\Bbbk$ 上有限型且具有仿射对角态射的代数叠之间的可表态射. 假定 $|\mathcal X|=\{x\}$ 与 $|\mathcal Y|=\{y\}$ 均为单点, 且 $f$ 诱导 $\mathcal X_0:=BG_x$ 与 $\mathcal Y_0:=BG_y$ 的同构. 令 $\mathfrak m_x\subseteq\mathcal O_X$, $\mathfrak m_y\subseteq\mathcal O_Y$ 为定义它们的理想层, $f_1:\mathcal X_1\to\mathcal Y_1$ 为一阶幂零加厚的诱导态射.
\begin{enumerate}[(1)]
\item 若 $f_1$ 是闭浸入, 则 $f$ 也是.
\item 若 $f_1$ 是闭浸入, 且有分次 $\mathcal O_{X_0}$-模的同构 $\bigoplus_{n\ge0}\mathfrak m_y^n/\mathfrak m_y^{n+1}\cong\bigoplus_{n\ge0}\mathfrak m_x^n/\mathfrak m_x^{n+1}$, 则 $f$ 是同构.
\end{enumerate}
\end{lemma}
\begin{proof}
取仿射概形的光滑展示 $V=\Spec B\to\mathcal Y$, 满足 $V\times_{\mathcal Y}\mathcal Y_0\cong\Spec\Bbbk$ (定理~\bookref{4.6.1}). 因 $\mathcal Y$ 仅有一个点, $B$ 是局部 Artin $\Bbbk$-代数. 基变换 $U=V\times_{\mathcal Y}\mathcal X$ 是代数空间, 因 $U_{\mathrm{red}}=V_{\mathrm{red}}$ 是点, 命题~\bookref{5.5.19} 说明 $U=\Spec A$, $A$ 为局部 Artin $\Bbbk$-代数. 故可假定 $f:\Spec A\to\Spec B$ 为局部 Artin 概形的态射. 对于 (1), 要证明 $B/\mathfrak m_B^2\to A/\mathfrak m_A^2$ 满射蕴含 $B\to A$ 满射. 先断言包含 $\mathfrak m_BA\to\mathfrak m_A$ 满射. 由 Nakayama 引理, 只需 $\mathfrak m_BA/\mathfrak m_A\mathfrak m_BA\to\mathfrak m_A/\mathfrak m_A^2$ 满射, 而假设说复合 $\mathfrak m_A/\mathfrak m_A^2\to\mathfrak m_BA/\mathfrak m_A\mathfrak m_BA\to\mathfrak m_A/\mathfrak m_A^2$ 满射. 因 $B/\mathfrak m_B\to A/\mathfrak m_BA=A/\mathfrak m_A$ 满射, 再用 Nakayama 引理, 得到 $B\to A$ 满射. 相关判据见 \cite[引理 II.7.4]{Har77}. 对于 (2), 因 $\dim_{\Bbbk}\mathfrak m_B^n/\mathfrak m_B^{n+1}=\dim_{\Bbbk}\mathfrak m_A^n/\mathfrak m_A^{n+1}$, 满射 $\mathfrak m_B^n/\mathfrak m_B^{n+1}\to\mathfrak m_A^n/\mathfrak m_A^{n+1}$ 为同构, 从而 $f$ 是同构.
\end{proof}

\subsection{局部结构定理的证明: 光滑情形}\label{src:7.7.3-subsection}
若 $x\in|\mathcal X|$ 是具有线性约化稳定子的光滑点, 将证明存在平展态射 $[T_{\mathcal X,x}/G_x]\leftarrow[\Spec A/G_x]\to\mathcal X$. 一般情形的证明不依赖光滑情形, 技术上复杂得多, 但思路相同. 因此建议先理解光滑情形.
\begin{proof}[定理~\bookref{7.7.1} 的证明: 光滑情形]
因 $\Bbbk$-点 $x$ 局部闭 (命题~\bookref{4.5.16}), 换为开子叠后可假定它闭. 令凝聚理想层 $\mathcal I$ 定义 $\mathcal X_0:=BG_x\hookrightarrow\mathcal X$, $\mathcal X_n$ 为 $\mathcal I^{n+1}$ 定义的第 $n$ 阶幂零加厚. Zariski 切空间 $T_{\mathcal X,x}$ 可与轨道的法空间 $(\mathcal I/\mathcal I^2)^\vee$ 等同, 视为 $G_x$-表示. 若 $\mathcal X=[X/G]$, $G$ 为光滑仿射代数群, 则对 $x$ 上的点 $\widetilde x\in X(\Bbbk)$, $T_{\mathcal X,x}$ 可与轨道的法空间 $T_{X,\widetilde x}/T_{G\widetilde x,\widetilde x}$ 等同.

定义 $\mathcal T:=[T_{\mathcal X,x}/G_x]$, 令 $\mathcal T_0=BG_x$ 为支撑在原点的闭子叠, $\mathcal T_n$ 为其各阶幂零加厚. 断言有相容同构 $\mathcal X_n\cong\mathcal T_n$. 因 $G_x$ 线性约化, $\mathcal X_0$ 上同调仿射. 主丛变形理论 (习题~\bookref{7.7.14}) 使我们将主 $G_x$-丛 $\Spec\Bbbk\to\mathcal X_0$ 逐阶延拓为 $\Spec A_n\to\mathcal X_n$, 得到同构 $\mathcal X_n\cong[\Spec A_n/G_x]$ 及仿射态射 $\mathcal X_n\to BG_x$. 从闭浸入 $\mathcal X_0\hookrightarrow\mathcal T$ 出发, 逐阶寻找提升
\begin{align*}
\begin{tikzcd}[ampersand replacement=\&]
\mathcal X_n\arrow[r]\arrow[d,hook]\&\mathcal T\arrow[d]\\
\mathcal X_{n+1}\arrow[r]\arrow[ur,dashed]\&BG_x,
\end{tikzcd}
\end{align*}
因为 $\mathcal T\to BG_x$ 光滑仿射 (命题~\bookref{7.7.12}). 诱导态射 $\mathcal X_1\to\mathcal T_1$ 是同构. 它是 $BG_x$ 由凝聚层 $\mathcal I/\mathcal I^2$ 作的变形之间的态射, 因此必为同构 (见引理~\bookref{C.1.7}). 事实上两者都可缩回到 $BG_x$, 故均为平凡变形. 引理~\bookref{7.7.16}(2) 说明各 $\mathcal X_n\to\mathcal T_n$ 为同构.

令 $\pi:\mathcal T\to T=T_{\mathcal X,x}//G_x$ 为到 GIT 商的态射. 纤维积 $\widehat{\mathcal T}:=\Spec\widehat{\mathcal O}_{T,\pi(0)}\times_T\mathcal T$ 是商叠 $[\Spec B/G]$, 其中 $B$ 在完备 Noether 局部 $\Bbbk$-代数 $B^G=\widehat{\mathcal O}_{T,\pi(0)}$ 上有限型. 所以 $\widehat{\mathcal T}$ 沿 $\mathcal T_0$ 凝聚完备 (定理~\bookref{7.6.13}), 凝聚 Tannaka 对偶~\bookref{7.6.9} 给出等价 $\operatorname{Mor}(\mathcal T,\mathcal X)\xrightarrow{\sim}\varprojlim\operatorname{Mor}(\mathcal T_n,\mathcal X)$. 因而 $\mathcal X_n\cong\mathcal T_n\hookrightarrow\mathcal X$ 延拓为 $\widehat{\mathcal T}\to\mathcal X$, 填充图
\begin{align*}
\begin{tikzcd}[ampersand replacement=\&,column sep=small]
\mathcal X_n\cong\mathcal T_n\arrow[r]\arrow[rrr,bend left]\&\widehat{\mathcal T}\arrow[r]\arrow[rr,dashed,bend left]\arrow[d]\&\mathcal T\arrow[d]\&\mathcal X\\
\&\Spec\widehat{\mathcal O}_{T,\pi(0)}\arrow[r]\&T.\&
\end{tikzcd}
\end{align*}
参数化态射同构类的函子
\begin{align*}
F:\Sch/T\to\Sets,\qquad(T'\to T)\longmapsto\{T'\times_T\mathcal T\to\mathcal X\}/\sim
\end{align*}
保持极限, 因 $\mathcal X$ 在 $\Bbbk$ 上有限型 (见习题~\bookref{4.3.33}). 态射 $\widehat{\mathcal T}\to\mathcal X$ 给出 $F$ 在 $\Spec\widehat{\mathcal O}_{T,\pi(0)}$ 上的元素. Artin 逼近~\bookref{B.5.18} 给出平展态射 $(U,u)\to(T,0)$, $U$ 为仿射概形且 $u$ 为 $\Bbbk$-点, 以及态射 $(U\times_T\mathcal T,(u,0))\to(\mathcal X,x)$, 它与 $(\widehat{\mathcal T},0)\to(\mathcal X,x)$ 一阶相同. 注意 $U\times_T\mathcal T\cong[\Spec A/G_x]$, $A$ 为有限生成 $\Bbbk$-代数且 $U=\Spec A^{G_x}$. 两边在相应点都光滑, 此态射诱导切空间与稳定子群的同构, 所以它在 $(u,0)$ 处平展. 将 $U$ 换为仿射开子概形后, 可安排该态射仿射 (命题~\bookref{7.7.17}) 且平展.
\end{proof}
\begin{proposition}\label{src:7.7.17}
设 $\mathcal X$ 是代数闭域上有限型且具有仿射对角态射的代数叠, $f:[\Spec A/G]\to\mathcal X$ 是有限型态射, $G$ 线性约化. 若 $w\in\Spec A$ 的 $G$-轨道闭, 且 $f$ 在 $w$ 处诱导稳定子群的同构, 则存在包含 $w$ 的 $G$-不变仿射开子概形 $U\subseteq\Spec A$, 使 $f|_{[U/G]}$ 仿射.
\end{proposition}
\begin{proof}
令 $\mathcal W=[\Spec A/G]$, $\pi:\mathcal W\to\Spec A^G$. 因 $f$ 在 $w$ 处诱导稳定子群同构, 将 $\Spec A^G$ 换为包含 $w$ 的像的仿射开集后, 可由上半连续性安排惯性叠 $I_{\mathcal W/\mathcal X}\to\mathcal W$ 拟有限. 取光滑展示 $V=\Spec B\to\mathcal X$, 考虑纤维积
\begin{align*}
\begin{tikzcd}[ampersand replacement=\&]
\mathcal W_V\arrow[r]\arrow[d]\&V=\Spec B\arrow[d]\\
{[\Spec A/G]}=\mathcal W\arrow[r,"{f}"']\&\mathcal X.
\end{tikzcd}
\end{align*}
因 $\mathcal X$ 的对角态射仿射, $\Spec B\to\mathcal X$ 仿射, 故 $\mathcal W_V$ 上同调仿射. 它的对角态射拟有限, 命题~\bookref{7.5.35} 说明 $\mathcal W_V\to V$ 分离, 下降说明 $\mathcal W\to\mathcal X$ 也分离. 因而相对惯性既固有又拟有限, 故有限. 在 $w$ 上的纤维平凡, 所以有开邻域 $\mathcal U$ 使相对惯性平凡. 再将 $\Spec A^G$ 换为仿射开集后, 可安排 $f$ 可表. 它也上同调仿射, 故 Serre 仿射性判据~\bookref{7.5.20} 说明它仿射.
\end{proof}

\subsection{等变 Artin 代数化}\label{src:7.7.4-subsection}
上述光滑性假设用于证明 $\mathcal T_n\cong\mathcal X_n$ 及 $U\times_T\mathcal T\to\mathcal X$ 平展. 更关键的是, 它说明 $\varprojlim\Gamma(\mathcal X_n,\mathcal O_{\mathcal X_n})$, 即极小泛变形空间的 $G_x$-不变量, 是有限生成 $\Bbbk$-代数的完备化 $\widehat{\mathcal O}_{T,0}$. 若 $x$ 不光滑, 很难直接建立这一点. 讨论 Artin 代数化~\bookref{C.6.8} 时也遇到类似问题: 若已知完备局部环 $R$ 是有限生成代数的完备化, 则 Artin 代数化是 Artin 逼近的简单推论 (见注记~\bookref{C.6.9}). 为绕过此问题, 一般证明将 $R$ 写成 $\widehat{\mathcal O}_{V,v}/I$, 其中 $V$ 是有限型 $\Bbbk$-概形, 如 $V=\mathbb A_{\Bbbk}^n$, $v=0$, 再用 Artin 逼近同时逼近 $R$ 上的给定对象与定义 $I$ 的方程. 这里采用类似策略, 但以 $G$-等变方式进行.

推广定义~\bookref{C.4.2} 的形式泛性: 若代数叠 $\widehat{\mathcal T}$ 有唯一闭点 $t$, 则 $\Sch$ 上预叠的态射 $\widehat\xi:\widehat{\mathcal T}\to\mathcal X$ 在 $t$ 处形式泛, 是指每个交换图
\begin{align*}
\begin{tikzcd}[ampersand replacement=\&]
\mathcal Z\arrow[r]\arrow[d,hook]\&\widehat{\mathcal T}\arrow[d,"{\widehat\xi}"]\\
\mathcal Z'\arrow[r]\arrow[ur,dashed]\&\mathcal X
\end{tikzcd}
\end{align*}
都有提升, 其中 $\mathcal Z\hookrightarrow\mathcal Z'$ 是具有仿射对角态射的 Noether 代数叠之间的闭浸入, $|\mathcal Z|=|\mathcal Z'|$ 为单点, 且 $\mathcal Z\to\widehat{\mathcal T}$ 的像为 $t$.
\begin{theorem}[等变 Artin 代数化]\label{src:7.7.18}
设 $\Bbbk$ 是代数闭域, $R$ 是剩余域为 $\Bbbk$ 的完备 Noether 局部 $\Bbbk$-代数. 设 $\widehat{\mathcal T}=[\Spec B/G]$ 是 $R=B^G$ 上有限型的代数叠, $G$ 线性约化. 假定唯一闭点 $t\in|\widehat{\mathcal T}|$ 的稳定子等于 $G$. 若 $\mathcal X$ 是 $\Sch/\Bbbk$ 上保持极限的预叠, $\eta:\widehat{\mathcal T}\to\mathcal X$ 在 $t$ 处形式泛, 则存在
\begin{enumerate}[(1)]
\item 有限型 $\Bbbk$-代数叠 $\mathcal W=[\Spec A/G]$ 及闭点 $w\in|\mathcal W|$;
\item 态射 $f:\mathcal W\to\mathcal X$ 与 $\varphi:\widehat{\mathcal T}\to\mathcal W$, 使图
\begin{align}\label{src:7.7.19}
\begin{tikzcd}[ampersand replacement=\&]
\widehat{\mathcal T}\arrow[d,"{\varphi}"']\arrow[dr,"{\eta}"]\&\\
\mathcal W\arrow[r,"{f}"']\&\mathcal X
\end{tikzcd}
\end{align}
中 $t,w$ 的第 $n$ 阶幂零加厚之间的诱导态射 $\varphi_n:\widehat{\mathcal T}_n\to\mathcal W_n$ 是同构, 并有相容 $2$-同构 $\eta_n\xrightarrow{\sim}f_n\circ\varphi_n$.
\end{enumerate}
若 $\mathcal X$ 还是有限型 $\Bbbk$-代数叠且对角态射仿射, 可安排 \bookeqref{7.7.19} 交换, $\varphi$ 诱导同构 $\widehat{\mathcal T}\to\widehat{\mathcal W}:=\mathcal W\times_W\Spec\widehat{\mathcal O}_{W,\pi(w)}$, 其中 $\pi:\mathcal W\to W=\Spec A^G$.
\end{theorem}
\begin{remark}\label{src:7.7.20}
取 $G$ 为平凡群, 就恢复经典 Artin 代数化~\bookref{C.6.8}.
\end{remark}
与 Artin 代数化的证明一样, 对选好的整数 $N$ 应用 Artin 逼近, 构造 $\mathcal W$, 使 $n\le N$ 时 $\mathcal W_n\cong\widehat{\mathcal T}_n$, 再用 Artin--Rees 引理~\bookref{B.5.4} 推出 $n>N$ 时也同构. 为控制 Artin--Rees 常数, 推广定义~\bookref{C.6.3}: 对 Noether 代数叠 $\mathcal X$, 由理想层 $\mathfrak m_x$ 定义的闭点 $x$ 与整数 $c\ge0$, 称凝聚层映射 $\varphi:E\to F$ 在 $x$ 处满足 $(\mathrm{AR})_c$, 若
\begin{align*}
\varphi(E)\cap\mathfrak m_x^nF\subseteq\varphi(\mathfrak m_x^{n-c}E),\qquad\forall n\ge c.
\end{align*}
当 $\mathcal X$ 为概形时, 即使将 $\{x\}$ 换为任意闭子概形, Artin--Rees 引理也给出充分大 $c$ 时成立. 光滑下降说明对代数叠, $c\gg0$ 时同样成立.
\begin{proof}
$\eta:\widehat{\mathcal T}\to\mathcal X$ 与 $\widehat{\mathcal T}\to BG$ 诱导 $\widehat{\mathcal T}\to\mathcal X\times BG$. 令 $\widehat T=\Spec R$ 为 $\widehat{\mathcal T}=[\Spec B/G]$ 的 GIT 商, $R=B^G$. $R$ 是有限生成 $\Bbbk$-子代数的余极限, $\mathcal X\times BG$ 保持极限, 极限方法 (\S\bookref{B.3}) 给出交换图
\begin{align*}
\begin{tikzcd}[ampersand replacement=\&]
\widehat{\mathcal T}\arrow[r]\arrow[d]\&\mathcal S\arrow[r]\arrow[d]\&\mathcal X\times BG\\
\widehat T\arrow[r]\&S.\&
\end{tikzcd}
\end{align*}
其中 $S=\Spec R'$ 是有限型仿射 $\Bbbk$-概形, $\mathcal S$ 是 $S$ 上有限型且具有仿射对角态射的代数叠, $\widehat{\mathcal T}=\widehat T\times_S\mathcal S$. 还可安排 $\mathcal S\to BG$ 仿射. 令 $\widetilde s\in|\mathcal S|$ 与 $s\in S$ 为 $t$ 的像. 添加 $R'$ 的生成元后, 可安排 $R'\to R\to R/\mathfrak m_R^2$ 满射. 完备 Nakayama 引理~\bookref{B.5.6}(3) 说明 $\widehat{\mathcal O}_{S,s}\to R$ 满射. 因而 $\widehat T\hookrightarrow\widehat S:=\Spec\widehat{\mathcal O}_{S,s}$ 与 $\widehat{\mathcal T}\hookrightarrow\mathcal S\times_S\widehat S$ 均为闭浸入.

取消解 $\mathcal O_{\widehat S}^{\oplus r}\to\mathcal O_{\widehat S}\to R$, 拉回到 $\mathcal S\times_S\widehat S$, 得到消解
\begin{align}\label{src:7.7.21}
\ker(\beta)\xrightarrow\alpha\mathcal O_{\mathcal S\times_S\widehat S}^{\oplus r}\xrightarrow\beta\mathcal O_{\mathcal S\times_S\widehat S}\longrightarrow\mathcal O_{\widehat{\mathcal T}}.
\end{align}
考虑函子 $F:\Sch/S\to\Sets$, 将 $S$-概形 $U$ 送到有限展示拟凝聚 $\mathcal O_{\mathcal S\times_SU}$-模的复形 $L\xrightarrow\alpha\mathcal O_{\mathcal S\times_SU}^{\oplus r}\xrightarrow\beta\mathcal O_{\mathcal S\times_SU}$ 的同构类集合. 极限论证说明 $F$ 保持极限. 复形 \bookeqref{7.7.21} 定义 $(\alpha,\beta)\in F(\widehat S)$, $\operatorname{coker}(\beta)=\mathcal O_{\widehat{\mathcal T}}$. 取整数 $N$ 使 $\alpha,\beta$ 在 $(\widetilde s,s)$ 处满足 $(\mathrm{AR})_N$. Artin 逼近~\bookref{B.5.18} 给出平展邻域 $(S',s')\to(S,s)$ 及 $(\alpha',\beta')\in F(S')$, 使两元素在 $F(\mathcal O_{S,s}/\mathfrak m_s^{N+1})$ 中相等. 令 $\mathcal W\hookrightarrow\mathcal S\times_SS'$ 为 $\operatorname{coker}(\beta')$ 定义的闭子叠, $w=(\widetilde s,s')$. 令 $S_n,S'_n,\widehat T_n$ 为 $S,S',\widehat T$ 在 $t$ 的像处的各阶幂零加厚, 则 $\widehat{\mathcal T}\times_{\widehat T}\widehat T_N$ 与 $\mathcal W\times_SS'_N$ 作为 $\mathcal S\times_SS_N$ 的闭子叠相等. 这对 $n\le N$ 给出 (1)--(2), 特别地有同构 $\varphi_N:\widehat{\mathcal T}_N\to\mathcal W_N$, 令其逆为 $\psi_N$.

由 $\eta$ 形式泛, 可对 $n\ge N$ 逐阶寻找相容提升
\begin{align*}
\begin{tikzcd}[ampersand replacement=\&]
\mathcal W_n\arrow[r,"{\psi_n}"]\arrow[d,hook]\&\widehat{\mathcal T}\arrow[d,"{\eta}"]\\
\mathcal W_{n+1}\arrow[r,"{f}"']\arrow[ur,dashed,"{\psi_{n+1}}"]\&\mathcal X.
\end{tikzcd}
\end{align*}
另一方面, 对复形 \bookeqref{7.7.21} 与 $(\alpha',\beta')$ 所定义复形的限制, 在 $\mathcal S\times_S\widehat S$ 上以 $c=N$ 应用引理~\bookref{C.6.4} (由光滑下降推广到叠), 得到分次 $\mathcal O_{BG}$-模的同构 $\operatorname{Gr}_{\mathfrak m_t}\mathcal O_{\widehat{\mathcal T}}\cong\operatorname{Gr}_{\mathfrak m_w}\mathcal O_{\mathcal W}$. 引理~\bookref{7.7.16} 说明所有 $\psi_n:\mathcal W_n\to\widehat{\mathcal T}_n$ 是同构. 因 $\widehat{\mathcal T}$ 凝聚完备 (定理~\bookref{7.6.13}), 凝聚 Tannaka 对偶~\bookref{7.6.9} 使逆映射 $\varphi_n=\psi_n^{-1}$ 有效化为 $\varphi:\widehat{\mathcal T}\to\mathcal W$, 完成 (1)--(2).

对于最后的断言, 若 $\mathcal X$ 代数, 再利用 $\widehat{\mathcal T}$ 与 $\mathcal W$ 的凝聚完备性, 对逆映射 $\psi_n$ 应用推论~\bookref{7.6.9}, 得到 $\psi:\widehat{\mathcal W}\to\widehat{\mathcal T}$, 它是 $\varphi$ 的逆. 因而 $\varphi:\widehat{\mathcal T}\to\widehat{\mathcal W}$ 是同构. 推论~\bookref{7.6.9} 的全忠实性给出延拓各 $\eta_n\xrightarrow{\sim}f_n\circ\varphi_n$ 的 $2$-同构 $\eta\to f\circ\varphi$, 故此同构是 $\mathcal X$ 上的态射.
\end{proof}

\subsection{局部结构定理的证明: 一般情形}\label{src:7.7.5-subsection}
\begin{proof}[定理~\bookref{7.7.1} 的证明]
可假定 $x$ 闭. 令 $\mathcal T=[T_{\mathcal X,x}/G_x]$, $\pi:\mathcal T\to T=T_{\mathcal X,x}//G_x$, $\widehat{\mathcal T}=\Spec\widehat{\mathcal O}_{T,\pi(0)}\times_T\mathcal T$. 令 $\mathcal T_0=BG_x$ 为支撑在原点的闭子叠, $\mathcal T_n$ 为其各阶幂零加厚. 将构造相容闭浸入 $\mathcal X_n\hookrightarrow\mathcal T_n$. 因 $G_x$ 线性约化, $\mathcal X_0=BG_x$ 上同调仿射. 用习题~\bookref{7.7.14} 变形主 $G_x$-丛 $\Spec\Bbbk\hookrightarrow\mathcal X_0$, 可逐阶构造同构 $\mathcal X_n\cong[\Spec A_n/G_x]$. 光滑仿射态射 $\widehat{\mathcal T}\to BG_x$ 的变形理论 (命题~\bookref{7.7.12}) 允许逐阶寻找提升
\begin{align*}
\begin{tikzcd}[ampersand replacement=\&]
\mathcal X_n\arrow[r]\arrow[d,hook]\&\widehat{\mathcal T}\arrow[d]\\
\mathcal X_{n+1}\arrow[r]\arrow[ur,dashed]\&BG_x.
\end{tikzcd}
\end{align*}
如光滑情形, $\mathcal X_1\to\mathcal T_1$ 是同构. 引理~\bookref{7.7.16}(1) 说明每个 $\mathcal X_n\to\mathcal T_n$ 都是闭浸入.

令 $I_n$ 定义 $\mathcal X_n\hookrightarrow\mathcal T_n$, 则 $\mathcal O_{\mathcal T_n}/I_n$ 是凝聚 $\mathcal O_{\mathcal T_n}$-模系统. 因 $\widehat{\mathcal T}$ 凝聚完备 (定理~\bookref{7.6.13}), 有凝聚理想层 $I\subseteq\mathcal O_{\widehat{\mathcal T}}$, 使满射 $\mathcal O_{\widehat{\mathcal T}}\to\mathcal O_{\widehat{\mathcal T}}/I$ 延拓各 $\mathcal O_{\mathcal T_n}\to\mathcal O_{\mathcal X_n}$. 它定义的闭浸入 $\widehat{\mathcal X}\hookrightarrow\widehat{\mathcal T}$ 延拓给定闭浸入, 给出实线箭头的交换图
\begin{align*}
\begin{tikzcd}[ampersand replacement=\&,column sep=small]
\mathcal X_n\arrow[r,hook]\arrow[d,hook]\&\widehat{\mathcal X}\arrow[d,hook]\arrow[r,dashed,"{\eta}"]\&\mathcal X\\
\mathcal T_n\arrow[r,hook]\&\widehat{\mathcal T}\arrow[r]\arrow[d]\&\mathcal T\arrow[d]\\
\&\Spec\widehat{\mathcal O}_{T,\pi(0)}\arrow[r]\&T.
\end{tikzcd}
\end{align*}
因 $\widehat{\mathcal X}$ 也凝聚完备, 凝聚 Tannaka 对偶~\bookref{7.6.9} 给出延拓此图的 $\eta:\widehat{\mathcal X}\to\mathcal X$. 因 $\widehat{\mathcal X}$ 与 $\widehat{\mathcal X}$ 有相同的幂零加厚, $\eta$ 在 $0$ 处形式泛. 以 $G=G_x$ 应用等变 Artin 代数化~\bookref{7.7.18}, 得到有限型 $\Bbbk$-代数叠 $\mathcal W=[\Spec B/G_x]$ 的态射 $f:\mathcal W\to\mathcal X$, 闭点 $w\in|\mathcal W|$, 以及 $\mathcal X$ 上的态射 $\varphi:\widehat{\mathcal X}\to\mathcal W$, 它诱导同构 $\widehat{\mathcal X}\to\mathcal W\times_W\Spec\widehat{\mathcal O}_{W,\pi(w)}$, $\pi:\mathcal W\to\Spec B^{G_x}$. 因 $f$ 诱导同构 $\mathcal W_n\to\mathcal X_n$, 它在 $w$ 处平展. 换为开子叠后可使 $f$ 处处平展. 命题~\bookref{7.7.17} 还使它仿射. 另见 \cite{AHR20,AHR19,AHHLR24}.
\end{proof}

\subsection{点处的凝聚完备化}\label{src:7.7.6-subsection}
称 $(\mathcal X,x)$ 为\newterm{完备局部叠}{complete local stack}, 若 $\mathcal X$ 是具有仿射稳定子的 Noether 代数叠, 有唯一闭点 $x$, 且沿剩余 gerbe $\mathcal G_x$ 凝聚完备. 定理~\bookref{7.6.13} 的商叠 $([\Spec A/G],x)$ 是重要例子, 如 $([\mathbb A^n/\mathbb G_m],0)$ 完备局部. Noether 代数叠 $\mathcal X$ 在点 $x$ 处的\newterm{凝聚完备化}{coherent completion} 是完备局部叠 $(\widehat{\mathcal X}_x,\widehat x)$, 配备态射 $\eta:(\widehat{\mathcal X}_x,\widehat x)\to(\mathcal X,x)$, 在 $\widehat x,x$ 的各阶无穷小邻域上诱导同构. 若 $\mathcal X$ 的对角态射仿射, 凝聚 Tannaka 对偶~\bookref{7.6.9} 说明 $(\widehat{\mathcal X}_x,\eta)$ 在唯一 $2$-同构意义下唯一.
\begin{theorem}\label{src:7.7.22}
设 $\mathcal X$ 是代数闭域 $\Bbbk$ 上有限型且具有仿射对角态射的代数叠. 对每个稳定子 $G_x$ 线性约化的 $x\in\mathcal X(\Bbbk)$, 凝聚完备化 $\widehat{\mathcal X}_x$ 存在, 且
\begin{enumerate}[(1)]
\item 它是商叠 $[\Spec B/G_x]$, 不变量环 $B^{G_x}$ 是有限型 $\Bbbk$-代数的完备化, $B^{G_x}\to B$ 有限型.
\item 若 $f:(\mathcal W,w)\to(\mathcal X,x)$ 平展, $\mathcal W=[\Spec A/G_x]$, $w$ 闭, 且 $f$ 在 $w$ 处诱导稳定子群的同构, 则 $\widehat{\mathcal X}_x\cong\mathcal W\times_W\Spec\widehat{\mathcal O}_{W,\pi(w)}$, 其中 $\pi:\mathcal W\to W=\Spec A^{G_x}$ 是到 GIT 商的态射.
\item 若 $\pi:\mathcal X\to X$ 是良模空间, 则 $\widehat{\mathcal X}_x=\mathcal X\times_X\Spec\widehat{\mathcal O}_{X,\pi(x)}$.
\end{enumerate}
\end{theorem}
\begin{proof}
局部结构定理~\bookref{7.7.1} 给出平展态射 $f:(\mathcal W,w)\to(\mathcal X,x)$, 其中 $\mathcal W=[\Spec A/G_x]$, $f$ 在闭点 $w$ 处诱导稳定子群的同构. 取 $\widehat{\mathcal X}_x=\mathcal W\times_W\Spec\widehat{\mathcal O}_{W,\pi(w)}$ 及 $B=A\otimes_{A^{G_x}}\widehat{A^{G_x}}$, 因 $\widehat{\mathcal X}_x=[\Spec B/G_x]$ 凝聚完备 (定理~\bookref{7.6.13}), 得到存在性及 (1), (2). (3) 由 (2) 及推论~\bookref{7.7.3} 得到.
\end{proof}
完备化决定有限型概形的平展局部结构 (推论~\bookref{B.5.21}), 这一事实可推广到叠.
\begin{theorem}\label{src:7.7.23}
设 $\mathcal X,\mathcal Y$ 是代数闭域 $\Bbbk$ 上有限型且具有仿射对角态射的代数叠, $x\in|\mathcal X|$, $y\in|\mathcal Y|$ 是具有线性约化稳定子群概形 $G_x,G_y$ 的 $\Bbbk$-点. 以下等价:
\begin{enumerate}[(1)]
\item 存在相容同构 $\mathcal X_n\to\mathcal Y_n$.
\item 存在同构 $\widehat{\mathcal X}_x\to\widehat{\mathcal Y}_y$.
\item 存在带 $G_x$-作用的仿射概形 $\Spec A$, $G_x$-不动点 $w\in\Spec A$, 以及平展态射的图
\begin{align*}
\begin{tikzcd}[ampersand replacement=\&]
\mathcal X\&{[\Spec A/G_x]}\arrow[l,"{f}"']\arrow[r,"{g}"]\&\mathcal Y,
\end{tikzcd}
\end{align*}
使 $f(w)=x$, $g(w)=y$, 且 $f,g$ 都在 $w$ 处诱导稳定子群的同构.
\end{enumerate}
若 $x,y$ 还光滑, 则这些条件等价于存在群概形同构 $G_x\to G_y$ 及在此同构下等变的切空间同构 $T_{\mathcal X,x}\to T_{\mathcal Y,y}$.
\end{theorem}
\begin{proof}
$(3)\Rightarrow(2)\Rightarrow(1)$ 直接得到; 凝聚 Tannaka 对偶~\bookref{7.6.9} 给出 $(1)\Rightarrow(2)$. 对于 $(2)\Rightarrow(3)$, 取局部结构定理~\bookref{7.7.1} 给出的平展邻域 $(\mathcal W=[\Spec A/G_x],w)\to(\mathcal X,x)$. 令 $\pi:\mathcal W\to W=\Spec A^{G_x}$ 为良模空间, 则 $\widehat{\mathcal X}_x\cong\mathcal W\times_W\Spec\widehat{\mathcal O}_{W,\pi(w)}$. 函子
\begin{align*}
F:\Sch/W\to\Sets,\qquad(T\to W)\longmapsto\operatorname{Mor}_{\Bbbk}(\mathcal W\times_WT,\mathcal Y)
\end{align*}
局部有限展示. 对 $F$ 及 $\alpha\in F(\Spec\widehat{\mathcal O}_{W,\pi(w)})$ 应用 Artin 逼近~\bookref{B.5.18}, 得到平展态射 $(W',w')\to(W,w)$ 及态射 $\varphi:\mathcal W':=\mathcal W\times_WW'\to\mathcal Y$, 使限制 $\varphi|_{\mathcal W'_1}:\mathcal W'_1\to\mathcal Y_1$ 是同构. 因 $\widehat{\mathcal W}_{w'}\cong\widehat{\mathcal X}_x\cong\widehat{\mathcal Y}_y$, 引理~\bookref{7.7.16} 说明 $\varphi$ 诱导同构 $\widehat{\mathcal W'}\to\widehat{\mathcal Y}$. 换为开邻域后得到平展态射 $(\mathcal W',w')\to(\mathcal Y,y)$. 最后断言来自 Luna 平展切片定理~\bookref{7.7.5}.
\end{proof}

\subsection{在等变几何中的应用}\label{src:7.7.7-subsection}
Sumihiro 环面作用定理~\bookref{B.1.34} 说, 带环面 $T$-作用的正规有限型 $\Bbbk$-概形的每个 $\Bbbk$-点都有 $T$-不变仿射开邻域. 不正规时未必如此, 如带 $\mathbb G_m$-作用的射影结点三次曲线 $C$ 的结点. 但总存在 $T$-等变仿射平展邻域, 如结点的 $\mathbb G_m$-等变平展邻域 $\Spec\Bbbk[x,y]/(xy)\to C$.
\begin{theorem}\label{src:7.7.24}
设 $X$ 是代数闭域 $\Bbbk$ 上局部有限型且具有仿射对角态射的代数空间, 带有仿射代数群 $G$ 的作用. 若 $x\in X(\Bbbk)$ 的稳定子线性约化, 则存在 $G$-等变平展邻域 $(\Spec A,u)\to(X,x)$, 在 $u$ 处诱导稳定子群的同构. 若 $G$ 是环面, 则每个 $x$ 都有这样的 $G$-不变平展邻域.
\end{theorem}
\begin{proof}
局部结构定理~\bookref{7.7.1} 给出平展邻域 $([\Spec A/G_x],u)\to([X/G],x)$, 其中 $u$ 闭且 $f$ 在 $u$ 处诱导稳定子群同构. 命题~\bookref{7.7.17} 允许将 $\Spec A$ 换为 $u$ 的 $G_x$-不变仿射开邻域, 使复合 $[\Spec A/G_x]\to[X/G]\to BG$ 仿射. 因而 $W:=[\Spec A/G_x]\times_{[X/G]}X$ 是仿射概形, $W\to X$ 是 $x$ 的 $G$-等变平展邻域. 若 $G$ 是环面, 它的任意子群都线性约化.
\end{proof}
